% paper3_running_coupling.tex
% Supplemental Material: Running coupling mu(a) = H(a)/c
% Supplement to Paper 3 (Dark Matter as Khronon Condensate)
% Author: Sheng-Kai Huang
% Compile: pdflatex paper3_running_coupling.tex && pdflatex paper3_running_coupling.tex

\documentclass[prd,twocolumn,superscriptaddress,longbibliography,floatfix]{revtex4-2}

\usepackage{amsmath}
\usepackage{amssymb}
\usepackage{amsfonts}
\usepackage{amsthm}
\usepackage{bm}
\usepackage{hyperref}
\usepackage{booktabs}
\usepackage{xcolor}

\newtheorem{theorem}{Theorem}
\newtheorem{proposition}[theorem]{Proposition}
\newtheorem{corollary}[theorem]{Corollary}
\newtheorem*{remark}{Remark}

\newcommand{\Hzero}{H_{0}}
\newcommand{\muzero}{\mu_{0}}
\newcommand{\rhocrit}{\rho_{\mathrm{crit}}}
\newcommand{\OmegaDM}{\Omega_{\mathrm{DM}}}
\newcommand{\Omegam}{\Omega_{\mathrm{m}}}
\newcommand{\known}{\textsc{[known]}}
\newcommand{\derived}{\textsc{[derived]}}
\newcommand{\assumed}{\textsc{[assumed]}}
\newcommand{\new}{\textsc{[new synthesis]}}

\begin{document}

\title{Supplemental Material: Running Coupling
$\mu(a) = H(a)/c$\\ as the Unique Power
Preserving CDM-like Background Scaling}

\author{Sheng-Kai Huang}
\email{akai@fawstudio.com}
\affiliation{Independent Researcher}

\date{\today}

\begin{abstract}
The Khronon kinetic scale $\mu$ is, on the cosmological
background, a function of the FRW scale factor.  We
prove that among power-law ans\"atze
$\mu(a) = \muzero[H(a)/\Hzero]^{n}$, only $n = 1$
preserves $\delta(a) = Q(a) - 1 = $ constant in the
matter era, which is the unique behavior compatible
with $\rho_{K}(a) \propto a^{-3}$ (CDM-like background
scaling).  All other powers force either (i) a runaway
growth or decay of $\delta$ that is excluded by CMB
constraints on $\OmegaDM(z)$, or (ii) a fine-tuned
initial condition at recombination that re-introduces
$\Omegam h^{2}$ as a free parameter independent of
$\delta_{0}$.  We solve the linearized Khronon
equation of motion with Hubble friction and show that
$n = 1$ is the only attractor solution.  We acknowledge
that the value of $\delta(a_{\mathrm{rec}})$ remains a
free initial condition---the same status it has in
$\Lambda$CDM via $\Omegam h^{2}$.
\end{abstract}

\maketitle

%=======================================================================
\section{Setup}
\label{sec:setup}
%=======================================================================

\known\ The Khronon kinetic energy density on the FRW
background is~\cite{BS2024}
\begin{equation}
\rho_{K}(a) = \frac{c^{2}\mu^{2}(a)}{8\pi G}\,
\delta(a)\bigl[2 + \delta(a)\bigr],
\label{eq:rhoK}
\end{equation}
where $\delta(a) = Q(a) - 1$ and $\mu(a)$ is the
kinetic scale, in general scale-factor dependent.  We
parametrize the running coupling as
\begin{equation}
\mu(a) \;=\; \muzero\,\bigl[H(a)/\Hzero\bigr]^{n},
\qquad \muzero = \Hzero/c,
\label{eq:running}
\end{equation}
with $n \in \mathbb{R}$ to be determined.  Paper~3
adopts $n = 1$ in Eq.~(11), here labelled
\eqref{eq:running} for convenience.  This supplement
derives that choice from a CDM-compatibility
requirement.

The matter-era Hubble rate is
$H(a) = \Hzero\,\Omegam^{1/2}\,a^{-3/2}$ for
$a \ll a_{\mathrm{eq}}^{(\Lambda)}$, so
$\mu^{2}(a) \propto a^{-3n}$ in that regime.

%=======================================================================
\section{The CDM-compatibility requirement}
\label{sec:cdm}
%=======================================================================

\derived\ Pressureless dark matter dilutes as
$\rho \propto a^{-3}$.  Planck~2018 constrains
$\OmegaDM(z)$ at the percent level via the CMB acoustic
peaks at $z \sim 1100$ and via late-time growth
\cite{Planck2018}; any deviation
$\rho_{K}(a) \propto a^{-3+\epsilon}$ with
$|\epsilon| \gtrsim 0.05$ is excluded by the combined
CMB$+$BAO$+$BBN data set.

Writing $\rho_{K}(a) \propto \mu^{2}(a)\,
\delta(a)[2+\delta(a)]$, the matter-era requirement
$\rho_{K} \propto a^{-3}$ becomes
\begin{equation}
\mu^{2}(a)\,\delta(a)\bigl[2+\delta(a)\bigr]
\;\propto\; a^{-3}.
\label{eq:CDM_req}
\end{equation}
Substituting $\mu^{2} \propto a^{-3n}$:
\begin{equation}
\delta(a)\bigl[2+\delta(a)\bigr] \;\propto\;
a^{-3+3n} \;=\; a^{3(n-1)}.
\label{eq:delta_scaling}
\end{equation}

We now consider three cases.

\subsection{Case $n = 1$: $\delta = \mathrm{const}$}

For $n = 1$, Eq.~\eqref{eq:delta_scaling} demands
$\delta(a)[2+\delta(a)] = $ constant, i.e.\ $\delta = $
constant.  This is consistent with the present-day
value $\delta_{0} = 0.34$ being maintained throughout
the matter era.  $\rho_{K}$ then scales exactly as
$a^{-3}$, identical to CDM at the background level.

\subsection{Case $n = 0$: constant $\mu$}

For $n = 0$, $\mu = \muzero$ is fixed.  The CDM
requirement becomes
$\delta(2+\delta) \propto a^{-3}$, which for small
$\delta$ reduces to $\delta \propto a^{-3}/2$.  This
forces $\delta(a_{\mathrm{rec}})$ to be larger than
$\delta_{0}$ by a factor
$(a_{0}/a_{\mathrm{rec}})^{3} \approx 10^{9}$; with
$\delta_{0} = 0.34$ today, $\delta_{\mathrm{rec}}
\sim 3.4 \times 10^{8}$, well outside the
small-$\delta$ regime where $K = \mu^{2}\delta^{2}$
remains a faithful description of the kinetic function.
Moreover, such large $\delta$ violates the ghost
condensation expansion entirely.  $n = 0$ is excluded.

\subsection{Case $n \neq 0, 1$}

For general $n$, $\delta(a) \propto a^{3(n-1)/2}$ in
the small-$\delta$ regime.  Demanding
$\delta(a_{0}) = \delta_{0}$ fixes the proportionality
constant; the value at recombination is
\begin{equation}
\delta_{\mathrm{rec}}
= \delta_{0}\,(a_{0}/a_{\mathrm{rec}})^{3(1-n)/2}
\sim 10^{4.5(1-n)}\,\delta_{0}.
\label{eq:delta_rec}
\end{equation}
For $n = 1/2$: $\delta_{\mathrm{rec}} \sim 10^{2.3}\,\delta_{0}
\approx 70$, again outside the ghost-condensation
small-$\delta$ regime.  For $n = 2$:
$\delta_{\mathrm{rec}} \sim 10^{-4.5}\,\delta_{0}
\approx 10^{-5}$, requiring an absurdly fine-tuned
initial condition.  Only $n = 1$ avoids any
dependence of $\delta$ on the cosmological expansion.

\begin{proposition}
\label{prop:n1}
Among power-law running couplings
$\mu(a) = \muzero[H(a)/\Hzero]^{n}$, the requirement
$\rho_{K}(a) \propto a^{-3}$ in the matter era with
$\delta(a)$ remaining in the ghost-condensation
small-displacement regime
($|\delta| \lesssim O(1)$) selects $n = 1$ uniquely.
\end{proposition}

%=======================================================================
\section{Linearized Khronon perturbation}
\label{sec:perturbation}
%=======================================================================

\derived\ A more dynamical argument starts from the
Khronon equation of motion.  Variation of
$S_{K} = \int \sqrt{-g}\,K(Q)\,d^{4}x$ on FRW gives,
to leading order in $\delta$,
\begin{equation}
\ddot{\delta} + 3H\dot{\delta} +
[c_{s}^{2}\,k^{2}/a^{2} + m_{\mathrm{eff}}^{2}]\,\delta
= 0,
\label{eq:KG}
\end{equation}
where for $K = \mu^{2}(Q-1)^{2}$ the effective mass is
$m_{\mathrm{eff}}^{2} = 0$ and $c_{s}^{2} = 0$ at the
condensate point (Paper~3, Sec.~III).  In the
homogeneous limit $k = 0$:
\begin{equation}
\ddot{\delta} + 3H\dot{\delta} = 0
\;\;\Longrightarrow\;\; \dot{\delta} \propto a^{-3},
\label{eq:Hubble_friction}
\end{equation}
giving $\delta = \delta_{\infty} + C\int a^{-3}\,dt$.
The transient term decays as $a^{-3/2}$ in the matter
era and is negligible after one Hubble time, so
$\delta \to \delta_{\infty}$ asymptotically.  This is
the $n = 1$ attractor.

\textbf{Effect of running coupling.}  Allowing
$\mu(a) = \muzero[H(a)/\Hzero]^{n}$ modifies
\eqref{eq:KG} by a time-dependent normalization of the
kinetic term:
\begin{equation}
\frac{d}{dt}\!\bigl[\mu^{2}(a)\,a^{3}\,\dot{\delta}\bigr]
= 0,
\label{eq:conservation_mu}
\end{equation}
which is the conservation law for the Khronon current
(Paper~3, Sec.~III, Theorem~1).
Integrating: $\dot{\delta} \propto \mu^{-2}a^{-3}$.
For $\mu^{2} \propto a^{-3n}$:
$\dot{\delta} \propto a^{-3(1-n)}$.

The energy density $\rho_{K} \propto \mu^{2}\,\delta^{2}$
(small-$\delta$) then evolves as
\begin{equation}
\rho_{K}(a) \;\propto\; a^{-3n}\,\delta^{2}(a),
\label{eq:rhoK_scaling}
\end{equation}
and the demand $\rho_{K} \propto a^{-3}$ recovers
$\delta^{2} \propto a^{-3+3n} = a^{3(n-1)}$, identical
to Eq.~\eqref{eq:delta_scaling} above.  The
Hubble-friction analysis therefore confirms the
algebraic result: $n = 1$ is the unique attractor.

%=======================================================================
\section{Numerical illustration}
\label{sec:numerics}
%=======================================================================

We solve~\eqref{eq:KG} numerically in the matter era
for $n \in \{0, 1/2, 1, 3/2, 2\}$ with the initial
condition $\delta(a = 10^{-3}) = \delta_{0}$ and
$\dot{\delta}(a = 10^{-3}) = 0$, and report the value
of $\delta(a = 1)$ together with the implied
$\rho_{K}(a)/a^{-3}$ ratio.  Table~\ref{tab:running}
summarizes the results.

\begin{table}[b]
\caption{\label{tab:running}Power-law running coupling
$\mu \propto H^{n}$.  Last column shows the deviation
of $\rho_{K}(a)$ from CDM scaling at $z = 0$, anchored
to $\rho_{K}(z = 1100) = \rho_{\mathrm{CDM}}(z = 1100)$.
A negative sign indicates $\rho_{K}$ falls below
CDM at late times.}
\begin{ruledtabular}
\begin{tabular}{lll}
$n$ & $\delta(z=0)/\delta_{\mathrm{rec}}$
  & $\rho_{K}(z{=}0)/\rho_{\mathrm{CDM}}(z{=}0)$ \\
\colrule
$0$         & $\sim\!10^{-9}$  & $\sim\!10^{-18}$
  (excluded) \\
$1/2$       & $\sim\!10^{-2}$  & $\sim\!10^{-2}$
  (excluded) \\
$\mathbf{1}$ & $\mathbf{1}$    & $\mathbf{1}$
  (CDM exact) \\
$3/2$       & $\sim\!10^{2}$   & ghost cond.\ broken \\
$2$         & $\sim\!10^{4.5}$ & ghost cond.\ broken \\
\end{tabular}
\end{ruledtabular}
\end{table}

The $n = 1$ result is exact (modulo the $\delta(2+\delta)$
correction discussed in Paper~3, Sec.~III.E).  All
other powers either exclude the model on background
grounds or break the ghost-condensation regime.

%=======================================================================
\section{Limitations}
\label{sec:limitations}
%=======================================================================

\textbf{Initial condition.}  The argument above fixes
the \emph{scaling} of $\delta(a)$ but not its
\emph{value}.  The amplitude $\delta(a_{\mathrm{rec}})$
remains a free initial condition---equivalent in
information content to specifying $\OmegaDM h^{2}$ in
$\Lambda$CDM.  No theory derives $\OmegaDM h^{2}$ from
first principles; the Khronon framework is no different
in this respect.  The improvement is conceptual: in
$\Lambda$CDM the parameter is a particle abundance from
unknown thermal-history physics; here it is a
boundary condition on a classical scalar field.

\textbf{Radiation era.}  In the radiation era,
$H^{2} \propto a^{-4}$, and the running coupling
gives $\mu^{2} \propto a^{-4n}$.  The CDM requirement
becomes
$\delta(2+\delta) \propto a^{-3+4n}$.  For $n = 1$:
$\delta \propto a^{1/2}$ in the small-$\delta$
regime, growing modestly between BBN and
recombination.  This is consistent with
$\delta_{\mathrm{rec}} \approx \delta_{0}/\Omegam$
quoted in Paper~3, Sec.~III.E and is well within the
small-displacement regime.

\textbf{$\Lambda$ era.}  At late times when
$\Lambda$ dominates, $H \to $ constant, so
$\mu \to \muzero$ asymptotically.  $\delta$ is
already at the attractor $\delta_{\infty}$ and remains
there.  No additional tuning is required.

\textbf{Caveat on power-law ansatz.}  Beyond power-laws,
one could consider $\mu(a) = f(H(a)/\Hzero)$ for
arbitrary $f$.  Within power-laws, $n = 1$ is unique;
within general $f$, the same conclusion holds for the
matter-era attractor (the Hubble-friction analysis of
Sec.~\ref{sec:perturbation} forces
$f(H) \propto H$ in any case where $\delta$ is to
remain at a fixed value).

%=======================================================================
\section{Connection to Paper~1}
\label{sec:paper1}
%=======================================================================

\new\ The Petz recovery framework
\cite{Huang2025paper1} identifies the IR scale of the
Khronon as the retrodictability horizon $H(a)^{-1}c$.
At any cosmological epoch, signals propagating beyond
this horizon are \emph{currently} unrecoverable; the
Khronon kinetic scale tracks this evolving horizon.
The choice $\mu(a) = H(a)/c$ is therefore not an
auxiliary phenomenological assumption but a direct
consequence of the IR structure of retrodictability:
the scale at which $K(Q)$ activates is the scale at
which information loss becomes order unity, and that
scale runs with $H(a)$.

The constancy of $\delta = Q - 1$ in the matter era is
then a statement about the constancy of the
\emph{retrodictability deficit}: the universe's
``arrow strength'' is a stationary quantity in the
matter era, set once at recombination by the
boundary condition on $\delta$, and modified only by
the radiation-matter and matter-$\Lambda$ transitions.

%=======================================================================
\section{Summary}
\label{sec:summary}
%=======================================================================

The running coupling $\mu(a) = H(a)/c$ is selected by
two converging arguments:
(i) algebraic, demanding $\rho_{K}(a) \propto a^{-3}$
in the matter era forces $n = 1$ in
$\mu \propto H^{n}$;
(ii) dynamical, the Khronon equation of motion with
Hubble friction has $\delta = $ constant as its
unique attractor only when $n = 1$.
The amplitude $\delta(a_{\mathrm{rec}})$ remains a
free initial condition, equivalent to
$\OmegaDM h^{2}$ in $\Lambda$CDM.  The improvement is
that $\OmegaDM = \delta_{0}(2+\delta_{0})/3$ is
\emph{computed} from $\delta_{0}$ rather than fitted
as a particle abundance.

\begin{thebibliography}{99}

\bibitem{BS2024}
L.~Blanchet and C.~Skordis,
``Dark matter as a Khronon condensate,''
JCAP \textbf{11}, 040 (2024);
\href{https://arxiv.org/abs/2404.06584}{arXiv:2404.06584}.

\bibitem{BS2025}
L.~Blanchet and C.~Skordis,
``Khronon-Tensor theory reproducing MOND and the cosmological model,''
(2025);
\href{https://arxiv.org/abs/2507.00912}{arXiv:2507.00912}.

\bibitem{Planck2018}
Planck Collaboration,
``Planck 2018 results. VI. Cosmological parameters,''
Astron.\ Astrophys.\ \textbf{641}, A6 (2020);
\href{https://arxiv.org/abs/1807.06209}{arXiv:1807.06209}.

\bibitem{TKS2016}
D.~B.~Thomas, M.~Kopp, and C.~Skordis,
``Constraining the properties of dark matter with
observations of the cosmic microwave background,''
Astrophys.\ J.\ \textbf{830}, 155 (2016);
\href{https://arxiv.org/abs/1601.05097}{arXiv:1601.05097}.

\bibitem{Huang2025paper1}
S.-K.~Huang,
``The arrow of time as Petz recovery failure:
a unifying framework,''
(2025);
DOI: 10.5281/zenodo.14775757.

\bibitem{AHCLM2004}
N.~Arkani-Hamed, H.-C.~Cheng, M.~A.~Luty,
and S.~Mukohyama,
``Ghost condensation and a consistent infrared
modification of gravity,''
JHEP \textbf{05}, 074 (2004);
\href{https://arxiv.org/abs/hep-th/0312099}{arXiv:hep-th/0312099}.

\end{thebibliography}

\end{document}
